% 张三慧第四版V2《电磁学、光学、量子物理》第18章 磁力
% 主教材扫描PDF物理页157--161，原编号18.1--18.30；非思考题。
% 图18.20--18.33等依赖源图，条件按主教材重述，审核权限与原版权均待确认。
\section*{第18章\quad 磁力}
\subsection*{章末习题18.1--18.30}

\paragraph{习题 18.1}
一粒子质量 \(0.5\,\mathrm g\)，
电量 \(+2.5\times10^{-8}\,\mathrm C\)，
初始水平速率 \(6.0\times10^4\,\mathrm{m/s}\)。
要使它仍保持水平方向运动，
应施加大小和方向怎样的匀强磁场？

\paragraph{习题 18.2}
图18.20中，电子以初速度
\(v_0=1.0\times10^7\,\mathrm{m/s}\)
在 A 点向上进入磁场，沿顺时针半圆从 A 到 C；
半圆直径 \(10\,\mathrm{cm}\)。
（1）求所需磁场大小和方向；
（2）求从A到C所需的时间。

\paragraph{习题 18.3}
动能 \(2.0\times10^3\,\mathrm{eV}\)
的质子以与匀强磁场
\(B=0.10\,\mathrm T\) 成 \(89^\circ\) 的
夹角入射，作螺旋运动。
求周期 \(T\)、螺距 \(h\) 和半径 \(r\)。

\paragraph{习题 18.4}
电视机显像管电子由
\(2.0\times10^4\,\mathrm V\) 电势差加速，
到屏幕距离 \(0.20\,\mathrm m\)，
受横向地磁场 \(0.5\times10^{-4}\,\mathrm T\)。
估计电子束偏转多少，是否明显？
忽略其他偏转磁场。

\paragraph{习题 18.5}
北京正负电子对撞机中电子
在周长 \(240\,\mathrm m\) 的储存环作轨道运动。
已知动量 \(p=1.49\times10^{-18}
\,\mathrm{kg\,m/s}\)，
求使电子偏转的磁感应强度。

\paragraph{习题 18.6}
蟹状星云电子的动量约
\(10^{-16}\,\mathrm{kg\,m/s}\)，
星云磁场约 \(10^{-8}\,\mathrm T\)。
求其回转半径。
若电子到中子星附近、
磁场约 \(10^8\,\mathrm T\)，
回转半径变为多少？
按题给动量不变作比较。

\paragraph{习题 18.7}
气泡室的磁场为 \(20\,\mathrm T\)，
高能质子垂直磁场飞过时留下
半径 \(3.5\,\mathrm m\) 的圆弧径迹。
求质子的动量和能量，
高能部分采用相对论能量动量关系。

\paragraph{习题 18.8}
来自太阳的电子以
\(0.80\times10^8\,\mathrm{m/s}\)
进入地球赤道上空高层范艾伦带，
当地磁场 \(4\times10^{-7}\,\mathrm T\)，
求绕磁场线圆周运动的轨道半径。
电子沿磁感线作缓慢螺旋运动到近地磁北极，
那里磁场 \(2\times10^{-5}\,\mathrm T\)，
若采用相同垂直速率估算，
半径又为多少？

\paragraph{习题 18.9}
回旋加速器两个D盒直径
\(75\,\mathrm{cm}\)，
匀强磁场 \(B=1.5\,\mathrm T\)。
加速氘核质量 \(3.34\times10^{-27}\,\mathrm{kg}\)，
电量 \(e\)。
（1）所用交流电压的频率；
（2）氘核离开加速器的最大动能（MeV）。

\paragraph{习题 18.10}
质谱仪构造见图18.22：
待测带电离子经过相互垂直的电场 \(E\)
和磁场 \(B\) 的速度选择器，
再进入匀强分析磁场 \(B'\)，
打在胶片上。
（1）证明距离 \(l\) 与质量关系
\[
 m=\frac{qBB'l}{2E}.
\]
（2）一次测量中氧16离子轨迹的偏转距离
\(l_1=29.20\,\mathrm{cm}\)，另一同位素离子
\(l_2=32.86\,\mathrm{cm}\)。
已知 \(m_1=16.00\,\mathrm u\)，
求后一氧同位素的质量。

\paragraph{习题 18.11}
铜片厚 \(d=1.0\,\mathrm{mm}\)，
在 \(B=1.5\,\mathrm T\) 的磁场中，
磁场垂直铜片表面。
铜片每立方厘米有
\(8.4\times10^{22}\) 个自由电子，
电子电量为 \(-e\)。
通电 \(I=200\,\mathrm A\) 时，
（1）求两侧 \(a,a'\) 间霍尔电势差；
（2）解释为何铜片宽度 \(b\)
对该电势差无影响。原图18.23定义正负位置。

\paragraph{习题 18.12}
图18.24中一半导体样品体积为 \(a\times b\times c\)，
沿 \(x\) 方向有电流，沿 \(z\) 方向施加磁场。
给 \(a=0.10\,\mathrm{cm}\)，
\(b=0.35\,\mathrm{cm}\)，
\(c=1.0\,\mathrm{cm}\)，
\(I=1.0\,\mathrm{mA}\)，
\(B=3000\,\mathrm G\)，
两侧电势差 \(U_{AA'}=6.55\,\mathrm{mV}\)。
（1）根据图示电势正负和右手规则判定
\(P\) 型还是 \(N\) 型半导体；
（2）求载流子的数密度。

\paragraph{习题 18.13}
掺砷硅片为 \(N\) 型半导体，
电子数密度 \(2\times10^{21}\,\mathrm{m^{-3}}\)，
电阻率 \(1.6\times10^{-2}\,\Omega\mathrm m\)。
用它作霍尔探头，尺寸
\(0.5\,\mathrm{cm}\times0.2\,\mathrm{cm}
\times0.005\,\mathrm{cm}\)。
长端接电压 \(1\,\mathrm V\)，
测得宽边霍尔电压 \(1.05\,\mathrm{mV}\)。
若最大表面与主向磁场垂直，求 \(B\)。

\paragraph{习题 18.14}
磁力用于输送液态金属等导电液体。
图18.25中长度为 \(l\) 的管道上
施加横向磁场 \(B\)，
同时通过沿磁场和管道均垂直的电流密度 \(J\)。
（1）证明液体两端压差
\(\Delta p=JlB\)；
（2）要产生
\(1.00\,\mathrm{atm}\) 压差，
\(B=1.50\,\mathrm T\)、
\(l=2.00\,\mathrm{cm}\)，
所需电流密度多大？

\paragraph{习题 18.15}
图18.26中利用霍尔效应测血流速度。
血管直径 \(2.0\,\mathrm{mm}\)，
磁场 \(0.080\,\mathrm T\)，
感应电势差 \(0.10\,\mathrm{mV}\)。
估算血液流速。

\paragraph{习题 18.16}
安培天平结构如图18.27：
一臂垂下有 \(n\) 匝矩形线圈，
位于匀强磁场中的下水平导线段长 \(l\)，
磁场 \(B\) 垂直纸面。
当线圈电流由 \(I\) 反向为 \(-I\) 后，
另一边须加质量 \(m\) 的砝码才能恢复平衡，
取 \(g=9.80\,\mathrm{m/s^2}\)。
（1）求 \(B\)；
（2）\(l=10.0\,\mathrm{cm},n=5,
 I=0.10\,\mathrm A,m=8.78\,\mathrm g\)，
数值计算。

\paragraph{习题 18.17}
细导线绕成矩形线圈，
长 \(20\,\mathrm{mm}\)、宽
\(10\,\mathrm{mm}\)，共1000匝，
置于均匀外磁场 \(B=1000\,\mathrm G\)，
电流 \(100\,\mathrm{mA}\)。
（1）图18.28(a)：磁场垂直线圈平面，
且线圈仅局部进入场区；
求各边受力、整圈净力和力矩。
（2）图18.28(b)：磁场平行于线圈平面，
求各边受力、净力和力矩，并验证
\(\boldsymbol\tau=\boldsymbol\mu\times\boldsymbol B\)。

\paragraph{习题 18.18}
方形线圈由外皮绝缘细铜线绕成，
共200匝，每边长 \(150\,\mathrm{mm}\)，
置于 \(B=4.0\,\mathrm T\) 外磁场，
电流 \(I=8.0\,\mathrm A\)。
求线圈磁矩大小及最大磁力矩。

\paragraph{习题 18.19}
均匀带电、半径 \(R\) 的介质圆盘
面电荷密度为 \(\sigma\)，
以角速率 \(\omega\) 绕垂直盘面中心轴旋转。
证明磁矩
\[
 \mu=\frac14\pi\omega\sigma R^4,
 \qquad
 \mu=\frac q{2m}L,
\]
\(q\) 为总电荷，\(m\) 为盘质量，
\(L\) 为同轴机械角动量。

\paragraph{习题 18.20}
中子虽总电荷为零却有磁矩。
按原书简化夸克模型：
一个电量 \(+2e/3\) 的上夸克和
两个电量 \(-e/3\) 的下夸克，
在半径 \(r\) 的同一圆周上以同一速率 \(v\)
运动；两个下夸克绕行方向相同，
上夸克绕行方向相反。
（1）求磁矩；
（2）若 \(r=1.20\times10^{-15}\,\mathrm m\)，
中子磁矩实验值
\(9.66\times10^{-27}\,\mathrm{A\,m^2}\)，
求此模型中的夸克速率。

\paragraph{习题 18.21}
电子自旋磁矩实验值约
\(0.928\times10^{-23}\,\mathrm{J/T}\)。
把电子类比为一个带电量 \(e\)、
半径 \(R\) 的经典圆电流环，
若电量绕行速度为 \(c\)，
求所需 \(R\)。
结合实验测得电子的特征半径
小于 \(10^{-18}\,\mathrm m\)，
讨论该经典自旋模型是否合理。

\paragraph{习题 18.22}
图18.29：无限长载流直导线
电流 \(I\)；同一平面内紧邻它放矩形线圈，
长为 \(l\)、宽为 \(b\)，
左侧距长导线 \(a\)。
线圈电流 \(I_1\)。
（1）求受磁力大小及方向；
（2）求受到的最大磁力矩。
取原图近侧线圈电流与长直导线同向。

\paragraph{习题 18.23}
无限长薄壁金属筒沿轴线方向有均匀
表面电流，面电流密度为 \(j\)
（单位 \(\mathrm{A/m}\)）。
求单位面积筒壁在自身磁场中受的
磁力大小及方向，说明如何使用
筒壁两侧的平均磁场。

\paragraph{习题 18.24}
将一片均匀分布着电流的无限大载流平面
放入外匀强磁场，
平面中的电流与该磁场垂直。
图18.30标出电流面左右两侧的
磁感应强度 \(\boldsymbol B_1,\boldsymbol B_2\)。
求该载流平面单位面积受磁力的大小及方向，
按磁压和两侧磁场的矢量关系给出结果。

\paragraph{习题 18.25}
两条无限长平行直导线
间距 \(5.0\,\mathrm{cm}\)，
各通过 \(30\,\mathrm A\) 电流。
求一条导线单位长度受到另一条的磁力。
若两线仅有电子沿导线漂移，电子速率
\(1.0\times10^{-3}\,\mathrm{m/s}\)，
解释两线的电力与磁力为何不能直接
用静止点电荷的力比代替。

\paragraph{习题 18.26}
图18.31：半径 \(R\) 的无限长
半圆柱面，面上沿轴向均匀分布总电流 \(I\)；
在其轴线上另有与它反向的无限长直电流 \(I\)。
（1）求轴上导线单位长度受到的磁力；
（2）若改用一根与半圆柱面\emph{同方向、
同大小电流}的无限直线导线取代半圆柱面，
使轴上线电流受到完全同样的力，
该等效导线应放在何处？

\paragraph{习题 18.27}
图18.32示电磁轨道炮：
子弹位于两根平行长导轨之间，
两轨各通过方向相反电流 \(I\)，
导轨圆柱半径为 \(r\)，
两轴距离为 \(a\)。
已由无限长导线模型给出子弹受磁力
\[
 F=\frac{\mu_0I^2}{2\pi}
 \ln\frac{a+r}{r}.
\]
导轨长 \(L=5.0\,\mathrm m\)，
\(a=1.2\,\mathrm{cm}\)，
\(r=6.7\,\mathrm{cm}\)，
子弹质量 \(m=317\,\mathrm g\)，
出口速度 \(4.2\,\mathrm{km/s}\)。
（1）平均加速度是重力加速度多少倍？
（2）需要多大电流？
（3）若总效率40%，需要多少千瓦电源功率？
忽略子弹从导轨末端起速这一假设的细节。

\textit{来源尺寸待核：}若 \(a\) 真为两圆柱\emph{轴距}，
则 \(a=1.2\,\mathrm{cm}<2r=13.4\,\mathrm{cm}\)，
两固体发生重叠，不能视为物理上分离的双导轨。
原始扫描图与单位需优先重新核对。

\paragraph{习题 18.28}
图18.33：通电流 \(I\) 的柔软导线
一端悬挂重物使张力为 \(T\)，
底部一段在垂直磁场 \(B\)
的区域形成半径 \(r\) 的圆弧。
（1）证明 \(r=T/(BI)\)；
（2）若把软导线类比为一条沿同圆轨迹运动、
电量为 \(-q\)、动量为 \(p=qT/I\) 的粒子，
说明磁力下圆周运动为何符合这一关系。

\paragraph{习题 18.29}
两质子某时刻相距 \(a\)：
质子1沿两质子连线、远离质子2
以速率 \(v_1\) 运动；
质子2沿垂直于连线方向
以速率 \(v_2\) 运动。
忽略 \(v_i/c\) 的高阶项，
求两质子此刻相互施加的总电磁力
大小与方向，讨论它们是否满足牛顿第三定律，
以及电磁场动量守恒。

\paragraph{习题 18.30}
钠原子激发态
\(3^2P_{3/2}\) 有效磁矩为
\(2.39\times10^{-23}\,\mathrm{J/T}\)。
在磁场中磁矩投影量子化，
磁矩与磁场夹角分别为
\(39.2^\circ,75^\circ,105^\circ,140.8^\circ\)。
若外场 \(B=2.0\,\mathrm T\)，
求四种可能的磁势能。

